% ============================================================
% PACOTES
% ============================================================

\usepackage{lmodern}
\usepackage[T1]{fontenc}
\usepackage[utf8]{inputenc}
\usepackage[brazil]{babel}
\usepackage{indentfirst}
\usepackage{nomencl}
\usepackage{color}
\usepackage{graphicx}
\usepackage{microtype}
\usepackage{booktabs}
\usepackage{csquotes}

% Pacotes adicionais
\usepackage{lipsum}

% Marcador em vermelho para referência ou número a completar antes da entrega
\newcommand{\pendente}[1]{\textcolor{red}{[#1]}}

% ============================================================
% CITAÇÕES
% ============================================================

\usepackage[
    brazilian,
    hyperpageref
]{backref}

\usepackage[
    alf,
    abnt-emphasize=bf,
    abnt-etal-cite=2,
    abnt-etal-list=0,
    abnt-etal-text=it,
    abnt-and-type=e,
    abnt-doi=expand,
    abnt-url-package=url,
    abnt-thesis-year=both
]{abntex2cite}
\usepackage{xurl} % deixa URL longa quebrar em qualquer ponto

% ============================================================
% FORMATAÇÃO DOS TÍTULOS
% ============================================================

\setsecheadstyle{\bfseries\MakeUppercase}

% ============================================================
% BACKREF
% ============================================================

\renewcommand{\backrefpagesname}{Citado na(s) página(s):~}

\renewcommand{\backref}{}

\renewcommand*{\backrefalt}[4]{
    \ifcase #1
        Nenhuma citação no texto.
    \or
        Citado na página #2.
    \else
        Citado #1 vezes nas páginas #2.
    \fi
}

% ============================================================
% APARÊNCIA DO PDF
% ============================================================

\definecolor{blue}{RGB}{41,5,195}

\makeatletter

\hypersetup{
    pdftitle={\@title},
    pdfauthor={\@author},
    pdfsubject={Modelo de artigo científico com abnTeX2 atualizado},
    pdfcreator={LaTeX with abnTeX2},
    pdfkeywords={
        abnt
    }{
        latex
    }{
        abntex
    }{
        abntex2
    }{
        artigo científico
    }{
        NBR 6022
    }{
        NBR 10520
    },
    colorlinks=true,
    linkcolor=blue,
    citecolor=blue,
    filecolor=magenta,
    urlcolor=blue,
    bookmarksdepth=4
}

\makeatother

% ============================================================
% ÍNDICE
% ============================================================

% \makeindex

% ============================================================
% MARGENS
% ============================================================

\setlrmarginsandblock{3cm}{2cm}{*}
\setulmarginsandblock{3cm}{2cm}{*}

\checkandfixthelayout

% ============================================================
% ESPAÇAMENTOS
% ============================================================

\setlength{\parindent}{1.25cm}
\setlength{\parskip}{0.2cm}

\SingleSpacing

% Para espaçamento 1,5:
% \OnehalfSpacing

\setlength{\ABNTEXcitacaorecuo}{4cm}